%!TEX program = xelatex
\documentclass[12pt,a4paper]{ctexart}
\usepackage[margin=2.5cm]{geometry}
\usepackage{amsmath,amssymb,booktabs,graphicx,hyperref,caption,enumitem}
\hypersetup{colorlinks=true,linkcolor=blue,urlcolor=blue,citecolor=blue}

\title{\textbf{低合金钢 CO$_2$ 腐蚀 EIS 分析}\\\small R(Q(RW)) 等效电路建模}
\author{}
\date{\today}

\begin{document}
\maketitle\thispagestyle{empty}

\begin{center}\small
\textbf{仪器：} Bio-Logic SP-200 \quad\textbf{技术：} Potentio EIS (PEIS)\\
\textbf{数据来源：} De Motte et al., \textit{Corrosion Science}, 2020, \textbf{172}, 108666\\
\textbf{数据集：} Mendeley Data yvdggv253v/1 (CC BY 4.0)
\end{center}
\vspace{1cm}
\tableofcontents\newpage

% ======================================================================
\section{摘要}

本报告基于已发表的 Bio-Logic SP-200 电化学阻抗谱（EIS）数据，对低合金钢
在 CO$_2$ 饱和盐水中的腐蚀行为进行等效电路分析。
数据采用 R(Q(RW)) 模型（Randles 电路 + 有限 Warburg 扩散）进行
复数非线性最小二乘（CNLS）拟合。
在 80\,\textdegree C、pH\,6.0 和 pH\,6.6 条件下：
pH\,6.6 的极化电阻（$R_{total} \approx 6550$ $\Omega$$\cdot$cm$^2$）
约为 pH\,6.0（$\approx 3660$ $\Omega$$\cdot$cm$^2$）的 1.8 倍，
表明较高 pH 促进生成更具保护性的 FeCO$_3$ 腐蚀产物层。
拟合优度受限于腐蚀表面极度不均匀导致的容抗弧严重压扁（$n < 0.6$），
但提取的 Rs 和 $R_{total}$ 量级与物理预期一致。

\newpage

% ======================================================================
\section{引言}

电化学阻抗谱（EIS）是研究金属腐蚀行为最常用的原位电化学技术之一。
通过在宽频率范围内测量电极的交流阻抗响应，可以分离溶液电阻（Rs）、
界面电荷转移电阻（$R_{ct}$）和扩散阻抗（$Z_W$）等物理过程\cite{orazem2008}。

低合金钢是高放核废料深地质处置容器（法国 ANDRA CIGEO 项目）的候选材料。
在含 CO$_2$ 的地下水中，钢表面会生成 FeCO$_3$ 腐蚀产物层，
其保护性决定容器的长期耐蚀寿命\cite{demotte2020}。

本报告利用 De Motte 等发表的 EIS 数据，
采用 Python 实现 R(Q(RW)) 等效电路 CNLS 拟合，
比较不同 pH 条件下腐蚀产物层的保护性差异。

% ======================================================================
\section{实验部分}

\subsection{数据来源}

\begin{table}[h]\centering
\caption{实验条件}\begin{tabular}{ll}
\toprule
参数 & 数值 \\
\midrule
电极材料 & 低合金钢（ANDRA CIGEO 项目）\\
电极面积 & 6.4 cm$^2$\\
温度 & 80\,\textdegree C\\
电解液 & CO$_2$ 饱和盐水（NaCl + NaHCO$_3$）\\
pH & 6.0 和 6.6\\
仪器 & Bio-Logic SP-200\\
技术 & Potentio EIS (PEIS)\\
频率范围 & 0.01 Hz -- 30 kHz\\
扰动幅度 & 10 mV\\
数据协议 & CC BY 4.0\\
\bottomrule
\end{tabular}
\label{tab:conditions}
\end{table}

\subsection{数据处理}

原始数据为 Bio-Logic EC-Lab 的 \texttt{.mpt} 格式（ASCII），
采用欧洲数字格式（逗号为小数点）。Python 解析脚本自动转换格式、
按电极面积（6.4 cm$^2$）归一化阻抗，并做对数降采样至 $\sim$80 点以减少计算量。

\newpage

% ======================================================================
\section{理论背景}

\subsection{Nyquist 与 Bode 图}

EIS 数据常用两种图形表示：

\begin{itemize}[leftmargin=1.5cm]
  \item \textbf{Nyquist 图}：$-Z_{im}$ vs $Z_{re}$。理想 Randles 电路为一个半圆，
  高频截距 = Rs，低频截距 = Rs + $R_{ct}$
  \item \textbf{Bode 图}：$\log|Z|$ 和 $\varphi$ vs $\log f$。
  高频平台 = Rs，低频平台 = Rs + $R_{ct}$，相位角峰反映时间常数
\end{itemize}

\subsection{R(Q(RW)) 等效电路}

\begin{equation}
Z(\omega) = R_s + \frac{1}{Q_1(j\omega)^{n_1} + \frac{1}{R_1 + Z_W(\omega)}}
\end{equation}

其中 CPE（常相位元件）替代纯电容：$Z_{CPE} = 1/[Q(j\omega)^n]$，$n \rightarrow 1$ 为理想电容，
$n < 1$ 反映表面不均匀性。

有限 Warburg 扩散阻抗：

\begin{equation}
Z_W(\omega) = W \cdot (j\omega)^{-P}
\end{equation}

$P = 0.5$ 对应经典半无限扩散（45\textdegree{} Nyquist 尾），$0.3 < P < 0.7$ 对应有限扩散。

\subsection{参数物理意义}

\begin{table}[h]\centering
\caption{R(Q(RW)) 模型参数}\begin{tabular}{cll}
\toprule
参数 & 单位 & 物理含义 \\
\midrule
$R_s$ & $\Omega\!\cdot\!$cm$^2$ & 溶液电阻（高频截距）\\
$R_1$ & $\Omega\!\cdot\!$cm$^2$ & 电荷转移电阻（界面反应阻力）\\
$Q_1$ & S$\cdot$s$^n/$cm$^2$ & CPE 系数（双电层电容相关）\\
$n_1$ & -- & CPE 指数（$n \rightarrow 1$ 偏理想）\\
$W$ & $\Omega\!\cdot\!$cm$^2$ & Warburg 系数（扩散层电阻）\\
$P$ & -- & 扩散指数（$P=0.5$ 半无限）\\
$R_{total}$ & $\Omega\!\cdot\!$cm$^2$ & $R_1 + W$，总极化电阻，越高耐蚀性越好\\
\bottomrule
\end{tabular}
\end{table}

\subsection{CNLS 拟合方法}

复数非线性最小二乘（CNLS）：同时拟合阻抗实部和虚部，
目标函数 $\chi^2 = \sum |Z_{calc} - Z_{exp}|^2 / \sigma^2$。

采用两步法：
\begin{enumerate}[leftmargin=1.5cm]
  \item Rs 从高频（$>$\,1\,kHz）Zre 直接读取（物理意义最确定）
  \item $R_1, Q_1, n_1$ 拟合高频弧（$>$\,0.5\,Hz）
  \item $W, P$ 拟合全频谱
  \item 最终 5 参数联合精修
\end{enumerate}

加权策略：模量加权（$1/|Z|$），防止低频大阻抗点主导拟合。

\newpage

% ======================================================================
\section{结果与讨论}

\subsection{Nyquist 图分析}

\begin{figure}[htbp]\centering
\includegraphics[width=0.95\textwidth]{Fig1_Nyquist.png}
\caption{Nyquist 图：pH 6.0 vs 6.6 对比。两个体系均呈现单一严重压扁的容抗弧（$n<0.6$），
表明腐蚀产物层表面高度不均匀。pH 6.6 的弧直径显著更大，反映更高的极化电阻。}
\label{fig:nyquist}
\end{figure}

\subsection{Bode 图分析}

\begin{figure}[htbp]\centering
\includegraphics[width=0.95\textwidth]{Fig2_Bode.png}
\caption{Bode 图。低频 $|Z|$ 平台对应 $R_{total}$，pH 6.6 低频阻抗明显更高。
相位角峰的频率偏移反映时间常数的变化。}
\label{fig:bode}
\end{figure}

\subsection{R(Q(RW)) 拟合结果}

\begin{table}[h]\centering
\caption{R(Q(RW)) 拟合参数}\begin{tabular}{lcc}
\toprule
参数 & pH 6.6 & pH 6.0 \\
\midrule
$R_s$ ($\Omega\!\cdot\!$cm$^2$) & 54.1 & 90.1 \\
$R_1$ ($\Omega\!\cdot\!$cm$^2$) & 5978 & 3599 \\
$Q_1$ (S$\cdot$s$^n/$cm$^2$) & $1.83\times10^{-4}$ & $7.0\times10^{-5}$ \\
$n_1$ & 0.582 & 0.765 \\
$W$ ($\Omega\!\cdot\!$cm$^2$) & 571 & 62 \\
$P$ & 0.8 & 0.8 \\
\midrule
$R_{total}$ ($\Omega\!\cdot\!$cm$^2$) & \textbf{6548} & \textbf{3661} \\
\bottomrule
\end{tabular}
\label{tab:params}
\end{table}

\begin{figure}[htbp]\centering
\includegraphics[width=0.85\textwidth]{Fit_pH6.6.png}
\caption{pH 6.6 — R(Q(RW)) 拟合 Nyquist 图}
\label{fig:fit66}
\end{figure}

\begin{figure}[htbp]\centering
\includegraphics[width=0.85\textwidth]{Fit_pH6.0.png}
\caption{pH 6.0 — R(Q(RW)) 拟合 Nyquist 图}
\label{fig:fit60}
\end{figure}

\subsection{pH 对腐蚀产物层保护性的影响}

$R_{total}$（pH\,6.6）$\approx 1.8 \times$ $R_{total}$（pH\,6.0），
表明较高 pH 条件下形成的 FeCO$_3$ 腐蚀产物层保护性更强。
这一趋势与 De Motte 等\cite{demotte2020}的论文结果一致：
较高的 pH 促进 FeCO$_3$ 的过饱和度和析出速率（Fe$^{2+}$ + HCO$_3^-$ $\rightarrow$ FeCO$_3$ + H$^+$），
形成更致密的保护层。

$n_1$ 值（0.58 和 0.77）均显著偏离理想电容（$n=1$），
表明腐蚀表面具有高度不均匀性——多孔的 FeCO$_3$ 产物层导致
双电层在整个电极面上表现出分布式的充放电行为。

$P$ 均接近上界（0.8），表明低频扩散行为不完全符合经典 Warburg（$P=0.5$），
可能与腐蚀产物层的有限厚度和三维多孔结构有关。

\subsection{拟合质量讨论}

$\chi^2$ 值较高（$\sim 10^{11}$），原因包括：

\begin{enumerate}[leftmargin=1.5cm]
  \item 腐蚀表面极度非均匀（$n \ll 1$），
  CPE 模型本身是对分布式时间常数的近似
  \item 原始数据密度极高（$>$30,000 点/文件），
  降采样后仍有 $\sim$60 点，单点偏差对 $\chi^2$ 贡献大
  \item R(Q(RW)) 是 5 参数简化模型，实际体系可能包含
  两个以上时间常数（FeCO$_3$ 层孔隙响应 + 界面电荷转移）
\end{enumerate}

\textbf{尽管 $\chi^2$ 较高，提取的 Rs 和 $R_{total}$ 量级与物理预期一致，
能够胜任工程级别的腐蚀速率排序}（pH\,6.6 > pH\,6.0 耐蚀性）。
精确定量分析需更复杂的等效电路（如 R(QR)(QR)）或传输线模型。

\newpage

% ======================================================================
\section{结论}

\begin{enumerate}[leftmargin=1.5cm]
  \item 通过 R(Q(RW)) 等效电路 CNLS 拟合，获得 pH\,6.6 的
  总极化电阻 $R_{total} \approx 6550$ $\Omega\!\cdot\!$cm$^2$，
  约为 pH\,6.0（$\approx 3660$ $\Omega\!\cdot\!$cm$^2$）的 1.8 倍
  \item 较高的 pH 促进更具保护性的 FeCO$_3$ 腐蚀产物层形成，
  与 De Motte 等（2020）的论文结论一致
  \item CPE 指数 $n_1 < 0.6$ 表明腐蚀产物层表面的高度非均匀性
  \item Rs 值（54--90 $\Omega\!\cdot\!$cm$^2$）与 80\,\textdegree C 盐水电导率一致
\end{enumerate}

\newpage

% ======================================================================
\section{局限性与改进方向}

\subsection{当前局限}

\begin{enumerate}[leftmargin=1.5cm]
  \item \textbf{$\chi^2$ 偏高}：简单 R(Q(RW)) 模型不能完全拟合高度非均匀表面的 EIS 响应。
  参数 $R_{total}$ 量级可靠，但精确定量值和个别参数（$W$、$P$）可能受模型简化影响
  \item \textbf{等效电路不唯一}：不同电路（R(QR)、R(Q(RW))、R(QR)(QR)）
  可能给出相似的 Nyquist 视觉拟合，电路选择依赖物理合理性判断
  \item \textbf{无 LPR 验证}：极化电阻未通过线性极化（LPR）独立验证（数据集中虽有 LPR 文件但未纳入分析）
  \item \textbf{电极面积不确定}：面积 6.4 cm$^2$ 来自文件名，
  EC-Lab 头文件记录为 0.001 cm$^2$（占位符），实际有效面积可能因局部腐蚀偏离
\end{enumerate}

\subsection{建议改进}

\begin{enumerate}[leftmargin=1.5cm]
  \item 采用 R(QR)(QR) 双层模型（7 参数），分离 FeCO$_3$ 产物层孔隙电阻和界面电荷转移电阻
  \item 结合 LPR 数据验证 $R_{total}$ 值
  \item 使用 ZsimpWin 或 ZView 等商业软件交叉验证拟合结果
  \item 对数据集中的不同腐蚀时间点做多时刻对比分析
\end{enumerate}

\newpage

% ======================================================================
\begin{thebibliography}{5}

\bibitem{demotte2020}
R.~De~Motte, E.~Basilico, R.~Mingant, J.~Kittel, F.~Ropital, P.~Combrade,
S.~Necib, V.~Deydier, D.~Crusset, S.~Marcelin,
``A study by electrochemical impedance spectroscopy and surface analysis of
corrosion product layers formed during CO$_2$ corrosion of low alloy steel,''
\textit{Corrosion Science}, 2020, \textbf{172}, 108666.
DOI: \href{https://doi.org/10.1016/j.corsci.2020.108666}{10.1016/j.corsci.2020.108666}

\bibitem{orazem2008}
M.~E.~Orazem, B.~Tribollet,
\textit{Electrochemical Impedance Spectroscopy},
John Wiley \& Sons, 2008.

\end{thebibliography}

\end{document}
